\section{Discriminant \texorpdfstring{\(-7\)}{-7}}\label{sec:minus7}

Put
\[
 Q_7(x,y)=x^2+xy+2y^2,\qquad
 P_7=\{Q_7(x,y)>0:\gcd(x,y)=1\},
\]
and let \(C_7\) contain \(1\) and the positive integers supported on
primes congruent to \(1,2,4\pmod7\). In particular, \(2\) is allowed.

\begin{theorem}\label{thm:minus7-exact}
The exact sumsets are
\[
 C_7+C_7=\{n\ge2:7\nmid n\}\setminus\{11,29\},
 \qquad
 P_7+P_7=\{N\ge2:49\nmid N\}\setminus\{7,77,203\}.
\]
For every \(N\ge2^{69}\) with \(49\nmid N\),
\[
       \#\{(a,b)\in P_7^2:a+b=N\}\ge N^{2/3}/20000.
\]
Pairs are ordered and summands are numerical values.
\end{theorem}

We write \(Q,P,C\) for these objects within this section.
The proof first produces a uniform sieve inequality on a quaternion
norm shell. Exact interval certificates connect its infinite range
to a complete finite sumset computation. A final descent removes
split prime factors, including the exceptional quotients \(11,29\).

\subsection{The value criterion and the ramified prime}

Let \(R=\Z[\omega]\), where \(\omega=(1+\sqrt{-7})/2\).
The Minkowski ideal-class bound is
\((2/\pi)\sqrt7<2\), so every ideal class has a representative of
norm one. Thus \(R\) has class number one. At an inert prime,
divisibility of a norm forces divisibility of both coordinates.
At a split prime, a primitive element can use only one of the two
conjugate prime ideals. At the ramified prime, its exponent as an
ideal divisor is at most one: the square of that prime ideal is
\((7)\). Conversely, assigning each entire split prime power to
one of its principal prime ideals, and choosing a ramified exponent
zero or one, gives a primitive generator. Therefore
\begin{equation}
                         P=C\mathbin{\dot\cup}7C.               \label{eq:m7-value}
\end{equation}
This is a statement about values; individual coordinate witnesses
used for the shell sieve need not be primitive at split primes.

Multiplication by \(\pi=\sqrt{-7}\) has matrix
\[
 M=\begin{pmatrix}-1&-4\\2&1\end{pmatrix},\qquad
 M^2=-7I,\quad \det M=7,\quad Q(Mv)=7Q(v).
\]
For a primitive vector of norm coprime to \(7\), the image is still
primitive. Indeed the determinant treats all other primes, while
both output coordinates vanish modulo \(7\) precisely when
\(2x+y\equiv0\pmod7\), which would imply \(7\mid Q(x,y)\).

Modulo \(7\), \(Q(x,y)=(x+4y)^2\). Since \(-1\) is not a square,
two nonzero squares cannot sum to zero. Every element of \(C\) is
a nonzero square modulo \(7\), and \eqref{eq:m7-value} gives
\begin{equation}
      7m\in P+P\quad\Longleftrightarrow\quad m\in C+C.           \label{eq:m7-scale}
\end{equation}
In particular \(49\mid N\) is impossible for \(N\in P+P\).
Every split prime preserves \(C\) under multiplication.

\subsection{The shell and its finite quotients}

Consider
\[
 B=(-7,-1)_{\Q},\qquad
 O=R\oplus Rj,\qquad j^2=-1,\quad ja=\bar a j .
\]
Its reduced norm is \(Q\oplus Q\).
The algebra has discriminant \(7\), and this order is maximal
\cite[Theorem 15.5.5 and Example 15.5.7]{Voight}.
Its units are \(\{\pm1,\pm j\}\): the two nonnegative component
norms sum to one, and the units of \(R\) are \(\pm1\).
The projective unit group \(\Gamma\) consequently has order two.
The projective mass formula gives total ideal-class mass
\((7-1)/12=1/2\) \cite[Theorem 25.1.1]{Voight}.
The principal class already contributes \(1/|\Gamma|=1/2\).
There is therefore one right-ideal class.

For \(7\nmid n\), every integral right ideal of norm \(n\) is
principal and has four generators. The all-ideal degree formula
\cite[Propositions 41.3.1 and 41.3.6]{Voight} gives
\begin{equation}
 r_O(n):=\#\{\alpha\in O:\operatorname{nrd}(\alpha)=n\}
                              =4\sigma(n).                    \label{eq:m7-shell}
\end{equation}
Likewise ideal factorization in \(R\) gives
\begin{equation}
 r_Q(t)=2\sum_{d\mid t}\chi_{-7}(d)
       \le2M(X)\quad(1\le t\le X),\qquad
 M(X)=\max_{1\le t\le X}
       \prod_{\substack{\ell\mid t\\\chi_{-7}(\ell)=1}}
                       (v_\ell(t)+1).                         \label{eq:m7-divisor}
\end{equation}

For an odd inert prime \(q\), reduction gives
\(G_q=\mathrm{PGL}_2(\mathbb F_q)\).
At \(7\), reduction modulo the radical gives
\[
                 G_7=\mathbb F_{49}^{\,*}/\mathbb F_7^{\,*},
\]
a cyclic group of order eight. Use the quotient
\(\Gamma\backslash\prod G_\ell\), for auxiliary levels \(7\),
\(q\), or \(q,7\), with uniform probability measure.
The image of \(\Gamma\) injects at every such level.
At odd \(q\), \(j\) cannot be scalar: its trace is zero and its
norm is one. At \(7\), \(j\notin\mathbb F_7\).
This argument includes \(q=3\).

For index \(n\) coprime to \(7\) and the added levels, the
all-ideal Hecke operator has degree \(\sigma(n)\) and satisfies
\begin{equation}
 (T_nh)(x_0)=\frac14
          \sum_{\operatorname{nrd}\alpha=n}h(\rho(\alpha)).
                                                                    \label{eq:m7-Hecke}
\end{equation}
Principalization initially gives inverse generators. Projectively
\(\alpha^{-1}=\bar\alpha\), and conjugation permutes the shell,
which proves the displayed orientation and normalization.
The operator is self-adjoint and commutes with the right finite
group action.

The global norm characters at level \(7\) are \(1,\chi_7\), where
\(\chi_7\) is the quadratic residue-field norm character.
At joint level \(q,7\) they are precisely
\begin{equation}
                        1,\ \chi_7,\ \chi_q,\ \chi_7\chi_q,     \label{eq:m7-chars}
\end{equation}
with \(\chi_q\) the determinant square character.
To check completeness, the maximal local unit norm maps are onto
away from the auxiliary levels. Reduction kernels and projective
central scalars leave only the two square classes at each auxiliary
prime. Globally, reduced norms are all positive rationals by
the Hasse--Schilling norm theorem
\cite[Main Theorem 14.7.4]{Voight}; rational square classes give
exactly \eqref{eq:m7-chars}. The strong-approximation step in
Lemma~\ref{lem:hecke-transfer} excludes other constituents having
a finite-dimensional full smooth local factor at a split place.
It is the full local factor, rather than a space of level-fixed
vectors, to which that argument applies.

Consequently Lemma~\ref{lem:hecke-transfer} gives, on the
orthogonal complement of the characters \eqref{eq:m7-chars},
\begin{equation}
                         \|T_n\|\le\tau(n)\sqrt n.             \label{eq:m7-Ramanujan}
\end{equation}
At every prime dividing \(n\), including \(2\), the quaternion
algebra is split and the local order is maximal. The all-ideal
prime-power eigenvalue is a sum of \(a+1\) monomials in two
roots of absolute value \(\sqrt\ell\), so the constant in
\eqref{eq:m7-Ramanujan} is one.

One must not remove every character of a finite right action.
For example \(\Gamma\backslash G_7\simeq C_4\) has two quartic
characters. The norm-two shell has eight pure-component points
reducing to \(\Gamma\) and four mixed points \(\pm1\pm j\)
reducing to \(h=[1+j]\). Here \(h\ne1\) and \(h^2=1\), because
\((1+j)^2=2j\). Thus
\(T_2=2I+R_h\), with spectrum \(3,3,1,1\).
The global norm characters have eigenvalue \(3\), whereas the
quartic characters have eigenvalue \(1\) and remain in the
complement in \eqref{eq:m7-Ramanujan}.

\subsection{A simultaneous ramified and inert filter}

For \(7\nmid n\), put
\[
 E(n)=\frac{\tau(n)\sqrt n}{\sigma(n)},\qquad
                     \gamma(n)=\frac{3-\chi_{-7}(n)}4 .
\]
Let \(G(n)\) count shell points for which both field norms are
\(7\)-units. These norms are necessarily positive. The two
component axes reduce to the single base point of
\(\Gamma\backslash G_7\). Its indicator \(\delta\) has mean
\(1/4\), coefficient \(1/4\) on \(\chi_7\), and complementary
squared norm \(1/8\). Apply \eqref{eq:m7-Hecke} to the correlation
of \(\delta\) with itself, and use \eqref{eq:m7-Ramanujan}.
The resulting lower bound is
\begin{equation}
                 \frac{G(n)}{4\sigma(n)}
                       \ge\gamma(n)-\tfrac12 E(n).             \label{eq:m7-good7}
\end{equation}
If \(\chi_{-7}(n)=-1\), an axis is impossible already modulo \(7\);
then \(G(n)=4\sigma(n)\) exactly.

For inert \(q\nmid n\), the invertible residues with one field
component zero form the normalizer \(H_q\) of a nonsplit torus.
It contains the unit image and has order \(2(q+1)\), hence fraction
\[
                            \mu_q=\frac{2}{q(q-1)}
\]
in \(G_q\). The determinant character has mean zero on \(H_q\),
because the torus norm is onto \(\mathbb F_q^*\).

Retain the ramified filter on the joint quotient. Let \(u\) require
the \(q\)-coordinate in \(H_q\) and the \(7\)-coordinate in
\(\Gamma_7=\{1,j\}\). Let \(v\) require the same \(q\)-condition and
the complementary \(7\)-condition. Both are right-invariant under
\(H_q\times\Gamma_7\), and the base orbit is the support of \(u\).
Consequently
\[
                     (T_nv)(x_0)
                       =\frac{\langle u,T_nv\rangle}{\mathbb E u}.
\]
This identity does not require \(v(x_0)=1\).
Writing \(a=\mu_q/4\), the coefficients on the four characters
in \eqref{eq:m7-chars} are respectively
\[
             (a,a,0,0)\quad\hbox{and}\quad(3a,-a,0,0).
\]
The complementary squared norms are \(a-2a^2\) and \(3a-10a^2\).
If \(b_q(n)\) counts points bad at \(q\) inside the \(7\)-unit shell,
Cauchy's inequality and \eqref{eq:m7-Ramanujan} therefore give
\begin{equation}
\begin{aligned}
 \frac{b_q(n)}{4\sigma(n)}
 &\le \gamma(n)\mu_q+
          \sqrt{(1-\mu_q/2)(3-5\mu_q/2)}\,E(n)\\
 &\le \gamma(n)\mu_q+\tfrac74 E(n).                            \end{aligned}\label{eq:m7-smallq}
\end{equation}
For \(\chi_{-7}(n)=-1\), omit level \(7\). The ordinary \(H_q\)
indicator has complementary squared norm \(\mu_q(1-\mu_q)\);
the error in \eqref{eq:m7-smallq} improves to \(E(n)\).

If \(q\mid n\), anisotropy forces both coordinate blocks to be
divisible by \(q\) as soon as either is. The spoiled count is zero
when \(v_q(n)=1\), and otherwise is exactly \(G(n/q^2)\).
Thus
\begin{equation}
 \frac{b_q(n)}{4\sigma(n)}
       \le\frac{\sigma(n/q^2)}{\sigma(n)}
           <q^{-2}\le\gamma(n)\mu_q.                           \label{eq:m7-targetprime}
\end{equation}
No level-\(q\) operator is applied to an index divisible by \(q\).

For large primes use geometry instead. Row integration gives
\[
 S_Q(X):=\#\{(x,y):0<Q(x,y)\le X\}
 \le \frac{2\pi}{\sqrt7}X+(2+4/\sqrt7)\sqrt X\le3X .
\]
For integers \(X\ge36\), the last inequality follows from
\(\pi<22/7\), \(\sqrt7>66/25\), and
\(50/21+58/99<3\). The finite certificate checks \(1\le X\le35\);
flooring extends the result to real \(X\).
A nonzero \(q\)-divisible block is \(q\) times a nonzero vector.
Choose its orientation, count it by \(3n/q^2\), and use
\eqref{eq:m7-divisor} for the other block. Hence
\begin{equation}
                         b_q(n)\le12M(n)n/q^2.                \label{eq:m7-largeq}
\end{equation}
Only \(q\le\sqrt n\) can occur.

Let \(\mathcal I\) be the inert primes and put
\[
 B=\sum_{q\in\mathcal I}\mu_q,\quad
 \Pi(z)=\#\{q\in\mathcal I:q\le z\},\quad
 T(z)=\sum_{\substack{q\in\mathcal I\\q>z}}q^{-2}.
\]
A union bound using \eqref{eq:m7-good7}--\eqref{eq:m7-largeq}
leaves a shell proportion at least
\begin{equation}
 \gamma(n)(1-B)
       -E(n)\{\tfrac12+\tfrac74\Pi(z)\}-3M(n)T(z).              \label{eq:m7-sieve}
\end{equation}
In the negative-character case the sharper sufficient inequality is
\begin{equation}
                  E(n)\Pi(z)+3M(n)T(z)<1-B.                  \label{eq:m7-negative}
\end{equation}
A surviving pair has positive norms in \(C\), so proves \(n\in C+C\).

\subsection{Complete arithmetic envelopes}

The inert primes belong to \(3,5,13\pmod {14}\). Counting those
progressions, then bounding each decreasing tail by its first term
and integral, gives
\begin{equation}
 \Pi(z)\le3z/14+13/14,\qquad T(z)\le3/(14z)+3/z^2.              \label{eq:m7-progression}
\end{equation}
For the count, \(13/14\) is the largest discrepancy within one
period; this is checked over the fourteen residue positions.

Here and below, numerical certificate statements use exact integers
or rational numbers. Set \(R_0=2^{19}\), \(D_0=2^{64}\).
Let \(a_r\) be the least integer greater than \(R_0\) in residue
\(r\pmod {14}\), for \(r=3,5,13\), and put
\[
 \overline B=
 \frac1{D_0}\sum_{\substack{q\in\mathcal I\\q\le R_0}}
            \left\lceil\frac{2D_0}{q(q-1)}\right\rceil
 +\frac{2(R_0+1)}{R_0}
       \sum_{r=3,5,13}\left(\frac1{a_r^2}+\frac1{14a_r}\right).
\]
Since \(\mu_q\le2(R_0+1)/(R_0q^2)\) above \(R_0\), this is an
upper bound for \(B\). Exact evaluation gives
\begin{equation}
                         B\le\overline B<47/100.               \label{eq:m7-B}
\end{equation}

The finite Euler calculations also give
\begin{equation}
 M(X)\le8X^{1/6},\qquad
 E(n)<\frac{91}{20}n^{-1/3}\quad(n\hbox{ inert-supported}),\qquad
 E(n)\le15n^{-1/3}\quad(7\nmid n).                              \label{eq:m7-Euler}
\end{equation}
We give the completeness bounds for these computations.
For \(M(X)^6/X\), each split prime contributes
\((a+1)^6/\ell^a\). For \(\ell\ge64\) it is at most one,
using \(a+1\le2^a\). Below \(64\), exponents through \(8\) suffice:
the subsequent ratios are at most \((10/9)^6/2<1\).
The exact product of the finite maxima is below \(8^6\).

For \(E(n)^6n^2\), the prime-power factor is
\begin{equation}
                  W_\ell(a)=
               \frac{(a+1)^6\ell^{5a}}{\sigma(\ell^a)^6}.       \label{eq:m7-W}
\end{equation}
Its next/current ratio is strictly less than
\(((a+2)/(a+1))^6/\ell\).
At inert primes below \(64\), caps \(4\) at \(3\), \(3\) at \(5\),
and \(1\) at \(13,17,19,31,41,47,59,61\) therefore suffice.
All larger primes have factor at most one. Their exact product
is less than \((91/20)^6\).
For all primes other than \(7\), use caps \(8,4,3,2\) at
\(2,3,5,11\), and \(1\) at the other primes below \(64\).
Their product is less than \(15^6\). This proves every bound
in \eqref{eq:m7-Euler} without an upper bound on \(n\).

For smaller ranges we use exact envelopes. A maximizer defining
\(M(U)\) uses nonincreasing exponents on an initial segment of the
smallest split primes: interchanging out-of-order exponents lowers
the integer cost and preserves the divisor product.
A recursion through \(U=2^{54}\), recording every initial prefix
and nonincreasing exponent profile of cost at most \(U\), has
\(2556\) states. Its completeness follows because the next
prime-prefix product exceeds \(U\), and no exponent exceeds
\(\lfloor\log_2U\rfloor\).

For \(E\), enumerate the \(5120\) tuples within the inert caps
above, retaining integer cost and the product of
\eqref{eq:m7-W}. Partition them by
\[
                 (v_3(n),v_5(n))\in\{0,1,\ge2\}^2.
\]
Deleting primes at least \(64\) and capping the other exponents
increases the normalized factor and reduces the cost, preserving
this category. Thus the cost-restricted maxima dominate all
inert-supported \(n\le U\). Take their upward sixth roots with
denominator \(10^6\), verified by sixth powers, and call the result
\(c_{a,b}(U)\). On \(L\le n\le U\) in category \((a,b)\),
\begin{equation}
                      E(n)\le
                  \frac{c_{a,b}(U)}{\lfloor L^{1/3}\rfloor}.    \label{eq:m7-localE}
\end{equation}
The global frontier is defined in the same way without categories.

\subsection{The finite interval certificate}

A sharper large-prime count is used when \(\chi_{-7}(n)=1\).
The ordered nonzero square residues summing to \(n\pmod7\)
are uniquely \(n/2,n/2\). Define
\[
 S_r(X)=\#\{(x,y):0<Q(x,y)\le X,\ Q(x,y)\equiv r\pmod7\},
 \quad F(X)=\max_{r=1,2,4}S_r(X).
\]
If one block is \(q\alpha\), then
\(Q(\alpha)\equiv n/(2q^2)\pmod7\).
The complete ordinary point table is computed through \(X_0=2^{18}\).
For \(0\le y\le\lfloor\sqrt{4X_0/7}\rfloor\), set
\(h=\lfloor\sqrt{4X_0-7y^2}\rfloor\) and enumerate
\[
          -\lfloor(h+y)/2\rfloor\le x\le\lfloor(h-y)/2\rfloor.
\]
Weight a row by two when \(y>0\) and by one at \(y=0\);
simultaneous sign reversal supplies the omitted rows.
This lists every signed vector. The individual coefficients are
also checked against \(2\sum_{d\mid t}\chi_{-7}(d)\).
The first seven positive \(F\)-values are \(2,4,4,6,6,6,6\).

Fix an interval \([L,U]\) and include the first \(j\) inert primes
as small primes, always including \(3,5\).
Every omitted \(q\le\sqrt U\) is required to satisfy
\(\lfloor U/q^2\rfloor\le X_0\).
For category \((a,b)\in\{0,1,2\}^2\), where \(2\) denotes
valuation at least two, put
\[
 D_{a,b}=\frac{3^a}{\sigma(3^a)}\frac{5^b}{\sigma(5^b)},\qquad
 d=\#\{p\in\{3,5\}:p\mid n\}.
\]
Since \(p^a/\sigma(p^a)\) decreases with \(a\),
\(n/\sigma(n)\le D_{a,b}\).
For \(\chi_{-7}(n)=1\), the large-prime error is at most
\begin{equation}
 \frac{M(U)D_{a,b}}L
       \sum_{\substack{q\in\mathcal I\ {\rm omitted}\\q\le\sqrt U}}
                              F(\lfloor U/q^2\rfloor).         \label{eq:m7-Ftail}
\end{equation}
The factor is \(M\), because the two orientations and the bound
\(2M(U)\) for the other block are divided by \(4\sigma(n)\).

A known divisor \(p=3\) or \(5\) needs no spectral error.
In \eqref{eq:m7-targetprime} its cost is
\(\delta_p=0\) at valuation one and \(p^{-2}\) above it.
Its baseline refund is \(\gamma\mu_p-\delta_p\).
For \(e\) given by \eqref{eq:m7-localE}, a sufficient
positive-character test is consequently
\begin{equation}
 e\{\tfrac12+\tfrac74(j-d)\}
       +\hbox{the term in \eqref{eq:m7-Ftail}}
 <
 \tfrac12(1-\overline B)+
           \sum_{\substack{p=3,5\\p\mid n}}(\mu_p/2-\delta_p).
                                                                    \label{eq:m7-finetest}
\end{equation}
For the negative character use small error \(e(j-d)\),
large error
\(3M(U)D_{a,b}\sum_{q\ {\rm omitted},q\le\sqrt U}q^{-2}\),
and baseline \(1-\overline B+\sum(\mu_p-\delta_p)\).
Reciprocal squares are rounded upward with denominator \(D_0\).

\begin{lemma}[Exact interval certificate]\label{lem:m7-intervals}
Every inert-supported \(n\ge2^{24}\), and every split prime
\(n\ge2^{24}\), belongs to \(C+C\).
\end{lemma}

\begin{proof}
Here is the complete finite part of the certificate.
Partition \([2^{24},2^{37})\) at all points
\(2^k+j2^{k-4}\), \(24\le k\le36\), \(0\le j\le16\),
and at every jump in the \(M\)- or category \(E\)-frontiers.
The resulting \(257\) adjacent integer intervals are checked to
have no gap. On each interval test both character signs and
all nine valuation categories by \eqref{eq:m7-finetest} and its
negative-character version. Also test the positive-character
boundary \(\tau(n)\le4\), with \(e=4/\lfloor\sqrt L\rfloor\)
and no category discounts. There are \(257\cdot19=4883\) cases.
For each case choose the allowed prime cut minimizing its rational
error. This can be found by comparing the cost of one more
spectral term with the decrease of the large-prime tail; those
decreases are nonincreasing. Direct enumeration of every allowable
cut gives the same finite test.

The certificate verifies a margin greater than \(1/64\) in every
case, every point-table query lies at most \(X_0\), and every
required prime lies below \(R_0\). For example the worst case is
\[
 L=17825792,\quad U=18874367,\quad \chi_{-7}=1,\quad a=b=0,
\]
where \(M(U)=112\), \(e=1907201/261000000\), and the last included
prime is \(73\). The error is
\(563281558829/2271744000000\), below its baseline by more than
\(1/64\).

For \(37\le k\le53\), put \(L=2^k\), \(U=2^{k+1}-1\),
\(d_k=M(U)\), and let \(e_k\) come from the global \(E\)-frontier.
The seventeen exact tests are
\begin{equation}
 e_k\{\tfrac12+\tfrac74\Pi(z_k)\}+3d_kT(z_k)
                      <(1-\overline B)/2-1/64.                \label{eq:m7-highbins}
\end{equation}
The prime sum through \(R_0\) and the three progression integrals
define the rational tail. The cut again minimizes the displayed
error. All seventeen cases pass. The certificate additionally
checks \(4\lfloor L^{1/3}\rfloor\le\lfloor\sqrt L\rfloor\).
Since its global normalized \(E\)-constant is at least one,
\(4/\sqrt n\le e_k\); these tests cover split primes as well.

For the infinite tail \(n\ge2^{54}\), take
\(z=(7/4)n^{1/4}\) in \eqref{eq:m7-sieve}.
For inert-supported \(n\), \eqref{eq:m7-Euler} bounds the error by
\[
 \frac{185799}{31360}n^{-1/12}
      +\left(\frac{17}{8}\frac{91}{20}+\frac{1152}{49}\right)n^{-1/3}.
\]
Using \((45/2)^{12}<2^{54}\), \(n^{1/3}\ge2^{18}\), this is
less than \(53/200\), with margin
\(9574319/6165626880>1/1024\).
The baseline is greater than \(53/200\) by \eqref{eq:m7-B}.
For \(\tau(n)\le4\), \(E(n)\le4n^{-1/2}\le(91/20)n^{-1/3}\),
so the same infinite tail applies. The three consecutive ranges
cover the asserted families. All finite tables and integer
roundings are reproduced by the submitted certificate program.
\end{proof}

\subsection{Finite coverage and exceptional descent}

For \(B_0=2^{24}\), the submitted computation constructs every
\(P\)-membership flag in two ways: complete primitive coordinate
enumeration and the factorization criterion \eqref{eq:m7-value}.
The full bytearrays agree. Deleting multiples of \(7\) gives the
complete \(C\)-flags. There are respectively \(2490663\) and
\(2161776\) flagged values.

For completeness, the coordinate enumeration may restrict to
\(y\ge0\) and \(x\ge-y/2\), since simultaneous negation and
conjugation preserve norm and gcd. The inequality
\((2x+y)^2+7y^2\le4B_0\) then gives finite exact bounds.
For each value flag array \(A\), form its binary integer
\(B_A=\sum_{j\le B_0}A(j)2^j\) and the truncated union
\[
       U_A=\bigvee_{\substack{1\le a\le5000\\A(a)=1}}(B_A\ll a).
\]
Here \(\vee\) is bitwise union and \(\ll\) is an exact integer
shift. Every set bit has actual positive summands.
The \(1050\) primitive anchors and \(903\) core anchors cover every
locally eligible target through \(B_0\), except respectively
\(\{7,77,203\}\) and \(\{11,29\}\).
Every omitted bit is tested by a complete positive-pair search
through half the target; none has a representation. The union
also has no locally forbidden bit. No summand above \(B_0\)
can contribute to a target in this range.
Membership arrays are compared before hashing; the stored hashes
are identifiers of these exact finite objects, not substitutes
for reconstructing and comparing them.

Suppose there were a least exception \(n\) to the asserted
\(C+C\) identity outside \(11,29\). Finite coverage makes
\(n>B_0\). If a split prime \(q\mid n\), multiplicative closure
and leastness imply \(n/q\in\{1,11,29\}\).
The first case makes \(n\) a split prime and is covered by
Lemma~\ref{lem:m7-intervals}. In the other cases,
\(n=11q\) or \(29q\). Division by \(11\) or \(29\), which are
themselves split primes, shows that \(q\) must be \(11\) or \(29\);
otherwise a smaller represented target lifts. The only products
left are already represented:
\[
                 121=8+113,\qquad
                 319=2+317,\qquad
                 841=32+809.
\]
All summands lie in \(C\). With no split prime factor, \(n\) is
inert-supported, again contradicting Lemma~\ref{lem:m7-intervals}.
This proves the core identity. The exclusions modulo \(7\)
were proved earlier, and
\[
                  C\cap[1,29]=
                  \{1,2,4,8,11,16,22,23,29\}
\]
checks the two exceptional nonrepresentations directly.

For \(7\nmid N\), the core identity supplies primitive values
except at \(11,29\); the patches \(11=4+7\) and \(29=1+28\)
use values in \(P\). If \(7\mid N\), \eqref{eq:m7-scale}
shows that failure occurs exactly at quotient \(1,11,29\),
hence at \(7,77,203\). Multiples of \(49\) are impossible.
This proves the two exact identities in Theorem~\ref{thm:minus7-exact}.

\subsection{The explicit lower bound}

For arbitrary \(7\nmid n\), use the last bound in
\eqref{eq:m7-Euler}, put \(z=n^{1/4}\), and use the weaker
\(\Pi(z)\le3z/14+3/2\). The error in \eqref{eq:m7-sieve} is at most
\[
                         \frac{603}{56}n^{-1/12}
                                  +\frac{951}{8}n^{-1/3}.
\]
For \(n\ge2^{66}\), \(n^{1/12}>45\), \(n^{1/3}\ge2^{22}\);
the error is at most \(281051653/1174405120\).
Its distance below \(53/200\) is
\(150828519/5872025600>1/40\).
Since \(4\sigma(n)\ge4n\), at least \(n/10\) coordinate points
survive. Each ordered value pair has at most
\((2M(n))^2\le256n^{1/3}\) coordinate witnesses by
\eqref{eq:m7-divisor}. Thus it yields at least
\(n^{2/3}/2560\) ordered \(C\)-pairs.

For \(N\ge2^{69}\), \(49\nmid N\), write \(N=7^en\) with
\(e=0\) or \(1\). Then \(n>2^{66}\).
Multiply both core values by \(7^e\), using \eqref{eq:m7-value}.
This injects the value pairs into \(P^2\).
As \(7^{2/3}<4\), the resulting lower bound is stronger than
\(N^{2/3}/20000\), as required.
